\section*{0.1 One-Slide Summary}
\addcontentsline{toc}{section}{0.1 One-Slide Summary}

\textbf{Problem:}  
Demand forecasting treats consumer behavior as exogenous and uncontrollable.  
No existing framework closes the loop between demand prediction and active demand intervention.  
Supply chains absorb demand volatility as cost rather than managing it as a controllable variable.

\textbf{Idea:}  
Treat incentives as a real-time control signal.  
Optimize \emph{timing}, not just magnitude, using a constrained Markov decision process (CMDP) that adapts to behavioral state, habit strength, and fairness constraints.

\textbf{Key Mechanism:}
\begin{itemize}
  \item Behavioral state includes recency, habit, time since last incentive, budget state, and fairness state.
  \item BED policy selects a micro-incentive (BED Tokens) at each decision epoch.
  \item BED Tokens map into existing reward units (points, credits, miles, discounts) through loyalty APIs.
  \item Budget and fairness constraints enforced at the policy level.
\end{itemize}

\textbf{Architecture:}
\begin{itemize}
  \item \emph{Prototype:} QR-based, serverless, minimal-integration system for field experiments.
  \item \emph{Production:} Direct integration with loyalty systems, POS data streams, and behavioral state estimators.
\end{itemize}

\textbf{Key Results Targets:}
\begin{itemize}
  \item Demand lift $\geq 5\%$ relative to control at statistically significant levels.
  \item Cost-per-unit-shifted $<\$0.50$, ensuring economic sustainability.
  \item Heterogeneous treatment effects across consumer segments and store types.
\end{itemize}

\textbf{Next Steps:}
\begin{itemize}
  \item Complete simulation validation.
  \item Launch MIT CTL–sponsored field trial (40–200 retail nodes, 4–8 week window).
  \item Target publication in top-tier operations research venues.
\end{itemize}
